\section{Objetivos del Proyecto}

\subsection{Objetivo General}

Diseñar el sistema de ventilación y presurización positiva del laboratorio de 320 m$^3$ de volumen efectivo, garantizando 12 renovaciones de aire por hora y evitando la infiltración de contaminantes desde zonas adyacentes.

\subsection{Objetivos Específicos}

\begin{enumerate}
	\item Determinar el caudal de aire de impulsión requerido para alcanzar las 12 renovaciones de aire por hora en el recinto.
	\item Calcular la potencia del ventilador necesaria para vencer las pérdidas de presión del sistema de impulsión y exfiltración.
	\item Dimensionar el área de impulsión del ventilador (configuración circular) y las rejillas de exfiltración rectangulares.
	\item Sustentar el número de renovaciones de aire con base en la normativa internacional aplicable (ASHRAE, WHO, NIH, OSHA, NFPA).
	\item Generar los datos de contorno para un modelo CFD que permita verificar la distribución de velocidades y el campo de presiones en el laboratorio.
\end{enumerate}
